\chapter{研究结论与展望}

本章主要对研究的核心成果进行梳理，同时也对未来的研究方向做一简单的预测。系统总结关键发现之后，进一步探究不同的教育模式、自主创业动机以及外部激励这三者怎样交互影响大学生创业意愿的产生过程，在理论方面评价研究成果的价值，探讨它对于加深创业教育领域学术认识的意义，在实践角度提出可行的政策建议，助力高校创新创业教育体系的优化升级，并且全方位剖析研究的不足之处，确定进一步发展的方向，给后续的研究给予重要参照。

\section{研究结论}

本文以自我决定理论和计划行为理论为基础，探讨了不同类型创新创业教育课程对学生创业意愿的影响机制。研究得出以下三点主要结论：

首先，三类创新创业教育课程均能有效激发学生的自主动机。理论型课程通过系统化的知识传授帮助学生建立创业认知框架，实践型课程借助体验式学习增强学生的创业能力感知，而融合型课程则将理论与实践有机结合，形成知识迁移与技能转化的正向循环。三类课程均从提高学生认知认同的角度出发，使学生在内在学习动机上得到加强，其中融合型课程对自主动机的激发效果最为突出。这表明，无论采用何种教学模式，提升学生对创业的价值认同都是激发其内在动机的关键。

其次，自主动机在课程类型与创业意愿之间发挥着完全中介作用。研究发现，三类课程对创业意愿的影响均不是直接的，而是完全通过自主动机这一心理机制来实现。这意味着创新创业教育的核心价值不在于单纯传授创业知识和技能，而在于能否真正调动学习者内在的兴趣爱好与价值认同。只有当学生从内心深处认同创业的意义和价值时，课程内容才能有效转化为创业意愿，否则即使完成了课程学习，也难以产生真正的创业动力。

最后，外部动机仅在融合型课程中对自主动机与创业意愿的关系产生负向调节作用。在融合型课程中，当学生主要出于获得学分、奖励等外部激励而参与学习时，自主动机对创业意愿的促进作用会被削弱，即外部激励对内在动机产生了\nolinebreak[3]“挤出效应”。而在理论型和实践型课程中，外部动机并未表现出显著的调节作用。这一发现说明，外部激励的负面影响在需要学生高度投入的融合型课程中更为明显，而在学习投入要求相对较低的课程类型中则较难显现。

\section{理论贡献}

本研究的理论贡献主要体现在以下三个方面：

首先，本研究根据创新创业教育模式的理论框架，将其分为三种主要类型即理论导向型、实践驱动型和融合整合型，用实证研究的方法来探究不同类型的创新教育对学生自主学习动机的作用机制以及不同的影响。研究结果表明这三种类型课程都会提高学生的自主学习意愿，融合型课程效果最好。它加深了有关高效教学模式如何促进学生内在学习动机的认识，也给创新创业教育领域的实践给予了重要的数据支撑和科学根据。

其次，从研究结果可知，自主动机具有很强的完全中介效应。各类课程体验对于创业意愿的影响都要经过自主动机这个关键的中介变量。加深了对创新创业教育机理的认识，这类教育并不直接提高受试者的创业意愿，而是在激活个人自主性动机之后起作用。该结论又把自我决定理论的使用领域和理论内涵扩大到创业教育领域，证明了自我决定理论的适用性。

最后，研究数据表明外部动机存在着负向调节的特点。实证检验的结果表明，在融合型课程环境之下，外部奖励会对个体自主动机与创业意向之间的正相关关系产生负向影响，该现象体现出教育实践中外在激励取代内驱力的一种情况。它既加深了自己决定理论对于外在驱动力作用机制的认识，又拓宽了它的运用范畴到创业教育当中去。

\section{实践启示}

基于上述研究结论，本研究对高校创新创业教育实践提出以下三点启示：

首先，高校应致力于培养学生自主学习的意识和能力，将调动内在动机作为创新创业教育的核心任务。研究表明，自主动机是课程影响创业意愿的关键中介变量，这意味着单纯传授创业知识和技能难以达到预期效果，教育的关键在于激发学生内在的求知欲望和价值认同。高校可通过邀请创业成功者分享亲身经历、举办创业理念专题讲座、设计个性化的实践活动等方式，帮助学生深刻理解创业的本质与意义，从而有效提升其创业意愿和创业准备程度。

其次，高校应将融合型课程作为创新创业课程体系建设的主要发展方向。研究结果表明，融合型课程对激发学生自主动机的效果最为突出，其作用在于将理论知识与实践操作有机结合，形成知行合一的良性循环。在课程设计过程中，应注重采取\nolinebreak[3]“项目驱动”与\nolinebreak[3]“情境仿真”相结合的教学方式，让学生在真实或模拟的创业情境中开展学习活动，从而最大限度地发挥融合型课程的教育效能。

最后，高校应警惕外部激励的\nolinebreak[3]“挤出效应”，减少对功利性评价手段的依赖。研究发现，在学分、奖学金等外部报酬驱动下参与课程的学生，其内在学习动机对创业意愿的促进作用会被削弱。因此，高校在构建课程体系时应弱化以分数和奖励为导向的评价方式，着力培养学生内在的兴趣驱动和价值认同感，避免陷入唯分数、唯实用主义的教学误区。针对三类课程的不同特点，还应采取差异化的教学设计：理论型课程侧重通过案例分析和典型事件讲解来提高学生参与积极性，实践型课程注重理论指导与实践操作的结合以防止形式化倾向，融合型课程则需不断优化理论与实践的衔接方式，减少外部激励对教学效果的干扰。

\section{研究不足与展望}

本研究存在以下不足，有待未来研究进一步完善，主要体现在样本采集、研究设计和变量测量三个方面。

首先，本研究的样本规模较为有限，尤其是实践型课程和融合型课程的有效样本量相对不足，这可能对统计结论的稳健性产生影响。尽管理论型课程获得了较为充分的数据支持，但实践型课程和融合型课程分别仅有53例和33例有效样本，样本量之间的差距使得包含调节效应的复杂模型检验结果可信度受到一定制约。融合型课程作为一种兼具理论性与实践性的创新教学模式，其作用机理仍有待进一步探究。后续研究应完善数据采集方案，在扩大总体样本基数的同时，重视对融合型课程参与者的数据收集，从而提高研究结论的普适性和科学价值。

其次，本研究采用横截面调查设计，虽然能够揭示变量之间的相关关系，但在因果推断方面存在固有局限性。自主动机与创业意愿之间可能存在双向交互影响，具有初步创业动机的学生在接受专业教育后，其自我驱动行为模式可能会发生显著改变。考虑到创新创业人才培育效果具有长期性特点，仅凭横截面数据难以揭示政策干预的动态路径与发展趋向。未来研究应采用纵向追踪设计，深入考察创新创业教育效果的长期演变机制。

最后，本研究在变量测量方面存在一定不足，特别是外部动机量表的信度有待提升。外部动机量表的Cronbach's α系数为0.815，虽达到心理测量学的基本标准，但与其他核心变量相比仍有差距，这可能加大测量偏差并影响调节效应检验的可信度。此外，本研究仅以自主动机和外部动机为理论基础，未能纳入创业自我效能感等其他重要变量，而创业自我效能感作为感知行为控制的重要体现，在创业决策中可能同时发挥中介和调节作用。未来研究应优化外部动机量表设计，提高测量的科学性与精确度，同时构建更复杂的多层模型，系统分析各心理变量之间的动态交互关系。此外，本研究仅选取天津某高校创业学院为数据采集点，样本同质性较强，虽有助于控制外生变量干扰，但也限制了结论的推广范围。后续研究应拓展调查的地域范围和学校类型，涵盖更多高等教育机构，以增强理论发现的一般性和广泛适用性。

\clearpage
